% frontmatter/palabras/mapa-del-curso.tex -- el año entero en una tabla.
% La escalera de sílabas de la columna "Qué sílabas" es la misma que
% comprueba tools/gen_palabras.py (ESCALERA): si cambia allí, cambia aquí.
\section*{Mapa del curso}

\begin{center}
\renewcommand{\arraystretch}{1.4}
\begin{tabularx}{\linewidth}{@{}c c c >{\raggedright\arraybackslash}p{27mm} >{\raggedright\arraybackslash}X@{}}
\toprule
\textbf{Trimestre} & \textbf{Días} & \textbf{Época} & \textbf{Qué se lee} & \textbf{Qué sílabas} \\
\midrule
1 & 1--65    & otoño     & 1 palabra al día  & Solo sílabas directas: \emph{ma, pe, lo, su}. Primero palabras de
dos sílabas (\emph{ma·má, pa·to}); desde la semana 6, también de tres
(\emph{o·to·ño}). \\
2 & 66--130  & invierno  & 2 palabras al día & Llegan las sílabas que terminan en consonante (\emph{sol, pan,
bu·fan·da}), la \emph{c} y la \emph{g} suaves (\emph{ce, ci, ge, gi}) y la
\emph{h}, que no suena. \\
3 & 131--195 & primavera & 3 palabras al día & Ya valen todas: dos consonantes juntas (\emph{li·bro, flor}),
\emph{ch, ll, rr, qu, gu} (\emph{le·che, pe·rro, que·so}) y dos vocales
juntas (\emph{a·bue·la, nie·ve}). Palabras de hasta cuatro sílabas. \\
4 & 196--260 & verano    & 1 frase al día    & Ya no hay palabras sueltas: una frase corta, de dos palabras
(\emph{Toby espera.}) a cuatro (\emph{Dani lee un cuento.}). Es el
puente hacia el cuaderno de frases. \\
\bottomrule
\end{tabularx}
\end{center}

\vspace{6mm}
Cada trimestre tiene 13 semanas de 5 días -- también durante las
vacaciones: una página por día de lunes a viernes, en Navidad y en
verano igual que en el colegio. Cada semana cuenta un pedacito de la
historia de Lucía, Dani y Toby, con los mismos temas, semana a semana,
que el cuaderno de frases: si en casa hay dos hermanos, cada uno puede
llevar el suyo y esa semana leen los dos sobre lo mismo.

\vspace{4mm}
\textbf{Lucía} y \textbf{Toby} son las dos únicas palabras que se leen
``de un golpe'', como el propio nombre, desde el primer día -- aunque
todavía no toquen sus sílabas (la \emph{cí} de Lu·cí·a, la \emph{y} de
To·by). Son los nombres de la protagonista y de su perro, y aparecen
todo el año.

\vspace{4mm}
Al final de cada trimestre hay una medalla, y al final del cuaderno,
un diploma y una \textbf{clave de respuestas} (para \lblAdivina,
\lblUne, \lblEncuentra\ y \lblSiNo) por si hace falta comprobar algo.
\newpage
